\documentclass[12pt,a4paper]{article}
\usepackage[utf8]{inputenc}
\usepackage[T2A]{fontenc}
\usepackage[russian,english]{babel}
\usepackage{amsmath,amssymb,amsfonts}
\usepackage{graphicx}
\usepackage{hyperref}

\title{GRA Chiral Nullification Equilibrium \ \medskip Равновесие хирального обнуления в GRA}
\author{qqewq}
\date{\today}

\begin{document}
\maketitle

\begin{abstract}
\selectlanguage{english}
We present a formal framework for the chiral nullification dynamics, where the nullification operator $\mathcal{N}$ maps a system state $(\pi, w, s)$ into a chiral pair $(R, L)$ while setting $w=s=0$. The reverse process generates new foam $\Phi$. Equilibrium is achieved when the hierarchical stability rank $N$ satisfies $\partial^N \Phi / \partial k^N = 0$. We provide mathematical derivations and practical applications in natural and social sciences.

\selectlanguage{russian}
Представлена формальная структура хиральной динамики обнуления, где оператор $\mathcal{N}$ переводит состояние системы $(\pi, w, s)$ в хиральную пару $(R, L)$, обнуляя $w$ и $s$. Обратный процесс генерирует новую пену $\Phi$. Равновесие достигается, когда ранг иерархической стабильности $N$ удовлетворяет $\partial^N \Phi / \partial k^N = 0$. Приведены математические выводы и практические применения в естественных и гуманитарных науках.
\end{abstract}

\section{Introduction / Введение}
The concept of "foam" $\Phi$ represents internal contradiction or entropy in a multi-agent or hierarchical system. The nullification operator $\mathcal{N}$ is designed to reduce $\Phi$ to zero by transforming the policy $\pi$ into a chiral pair $(R, L)$, which are mirror images under reflection operator $Z$. However, the system never stays at zero foam; the chiral symmetry breaking generates new foam, creating a dynamic equilibrium.

\section{Mathematical Framework / Математическая структура}

\subsection{Foam Definition / Определение пены}
Let $\pi$ be a policy vector, $w$ weights, and $s$ state vector. Foam is defined as:
\begin{equation}
\Phi(\pi, w, s) = \|\pi - w\| + \|w - s\| + \|s - \pi\|.
\end{equation}
This measures the total pairwise inconsistency.

\subsection{Chiral Mapping / Хиральное отображение}
The reflection operator $Z$ acts component-wise: $(Z v)_i = (-1)^i v_i$ (for simplicity). Then:
\begin{equation}
R = \pi, \quad L = Z\pi.
\end{equation}
The pair $(R, L)$ is chiral if $\pi \neq Z\pi$ under any allowed transformation.

\subsection{Nullification Operator / Оператор обнуления}
\begin{equation}
\mathcal{N}(\pi, w, s) = \left( \begin{pmatrix} R \\ L \end{pmatrix}, 0, 0 \right).
\end{equation}
This simultaneously eliminates $w$ and $s$ (weights and state) and embeds the policy into a higher-dimensional chiral space.

\subsection{Hierarchical Stability / Иерархическая стабильность}
For a hierarchy of levels $k$, we define the rank $N$ as the smallest integer such that the $N$-th derivative of $\Phi$ with respect to $k$ vanishes:
\begin{equation}
N = \min\{ n \mid \partial^n \Phi / \partial k^n = 0 \}.
\end{equation}
At this point, the system switches from particle-like stability to wave-like behaviour.

\subsection{Wave Nullification / Волновая обнулёнка}
The wave function is given by:
\begin{equation}
\psi(k) = e^{-\lambda \Phi(k)},
\end{equation}
which describes the probability amplitude of the system being in a particular foam state.

\subsection{Equilibrium Condition / Условие равновесия}
Equilibrium between nullification and foam generation occurs when the system oscillates around $\Phi=0$ with amplitude governed by $N$. The balance is expressed as:
\begin{equation}
\frac{d\Phi}{dt} = -\alpha \Phi + \beta \psi(k),
\end{equation}
where $\alpha$ is the nullification rate and $\beta$ the generation rate. Steady state $\Phi_{eq} = \beta/\alpha \, e^{-\lambda \Phi_{eq}}$ (implicit).

\section{Applications / Применения}
\subsection{Natural Sciences / Естественные науки}
- Plasma confinement: minimizing foam corresponds to stabilizing instabilities.
- Superconductor search: multiverse levels filter contradictory models.

\subsection{Humanities / Гуманитарные науки}
- Collective intelligence: swarm subject formation through zero-foam agents.
- Decision-making: conflict resolution by nullifying contradictory objectives.

\section{Conclusion / Заключение}
The GRA chiral nullification equilibrium provides a universal framework for describing systems that alternate between coherence (zero foam) and exploration (foam generation). The interplay between the nullification operator and hierarchical stability rank $N$ offers a quantitative tool for modelling stability in complex adaptive systems.

\section*{References / Ссылки}
\begin{enumerate}
\item GRA-Core-new-Unified-Hierarchical-Stability-Library (github.com/qqewq)
\item GRA-Multiverse-Final (github.com/qqewq)
\item GRA-Chiral-Nullification-Math (github.com/qqewq)
\end{enumerate}

\end{document}
